\section{The trust model}
\label{sec:trust}

\Telperion's design rests on a single principle: \textbf{the generator is
untrusted, and the \Lean{} kernel is the sole trusted component.} A defective
certificate manifests as a compile failure, never as a false theorem. Everything
else in the tool---the search heuristics, the SDP rounding, the emitters, even the
exact-arithmetic self-checks---exists to make the tool \emph{useful}, not to make
its output \emph{true}. Truth is established by, and only by, the kernel
re-deriving the theorem from the emitted certificate.

This has a design corollary worth stating: because nothing the generator does is
trusted, the generator is free to be large and heuristic while the trusted surface
stays small. The engine is a few thousand lines of dependency-light Python (it
uses only \texttt{sympy}); a referee who wishes to can audit the engine, but need
not---the audit that matters is the kernel's, and it happens on every build. This
is the same discipline that underlies certificate-checking in optimization
(e.g.\ VIPR for mixed-integer programs~\cite{vipr}) and the ``de~Bruijn
criterion'' in interactive theorem proving~\cite{debruijn}: a small trusted checker
re-validates the output of an arbitrarily complex, untrusted producer.

\paragraph{The self-checks are a convenience, not the guarantee.} Before emitting,
\Telperion{} verifies the certificate numerically in exact arithmetic
(\texttt{fractions.Fraction} and \texttt{sympy.Rational}; no floating point on any
certificate path). This catches mistakes cheaply---before a continuous-integration
round-trip---but establishes nothing about truth. It is important to be precise
here, because it is a common category error: an exact-arithmetic check that ``the
certificate reproduces the target'' is a check of the \emph{generator's} internal
consistency. Only the kernel's re-derivation is evidence about the theorem.

\paragraph{The one thing the kernel cannot catch: vacuity.} The kernel rejects a
\emph{false} theorem. It happily accepts a \emph{true-but-vacuous} one: a reflexive
$X = X$ or $0 \le 0$ compiles green while proving nothing about the certificate,
because the defect lives in the \emph{statement}, not the proof. \Telperion{}
therefore turns the tool on its own output. The emit step refuses a reflexive
emitted statement, and the identity-style emitters additionally require the
certificate to be \emph{load-bearing}: a corrupted certificate must break the
claim, or emission is refused. A family that deliberately emits a reference
identity or a degenerate boundary case opts out explicitly, declaring that intent.
Two further local gates run before any build is attempted: a structural lint
(unfilled template holes, unbalanced delimiters, duplicate names) and a
\emph{soundness lint} that refuses the ``green build $\ne$ proved'' classes---a
\code{sorry} or \code{admit}, a smuggled \code{axiom}, an empty \code{:= by}, a
missing type ascription, a \code{Prop := True} stub.

\paragraph{Provenance and drift.} Emitted files are stamped with the tool version
and a SHA-256 hash over a canonical serialization of every instance's expressions,
the \Lean{} profile, the proof templates, and the emitters' own code (timestamps
excluded). A regeneration diff detects any drift byte-for-byte, so a change to
emission logic can never ship under a stale hash, and a hand-edit to a generated
file is flagged. There is no API path from a family definition to \Lean{} text that
bypasses certification, and emission is refused without both a certified-family
witness and a green validation report.

\paragraph{What kernel acceptance still does not give you.} Two gaps remain, and
they motivate \S\ref{sec:indep}. First, kernel acceptance rests on a \emph{single}
implementation of the kernel (and the elaborator that produced the proof term);
an implementation bug, however unlikely, is a common-mode failure. Second, the
non-vacuity guards above are \Telperion{} checking its \emph{own} output for a
narrow syntactic class of degeneracy; they are not an independent confirmation
that the emitted proof proves the statement a third party \emph{intended}. Both
gaps are closed---to the extent any finite process can---by handing the proof to
an independent judge that re-derives the intended statement's identity, restricts
the admissible axioms, and replays the proof through a second, independently
written kernel. That is the subject of \S\ref{sec:indep}; we first describe how
certificates become \Lean{} at all.
